\chapter{DHEA-Sulfate (Dehydroepiandrosterone Sulfate)}

\section*{What Is This Test?}
Dehydroepiandrosterone sulfate (DHEA-S) is the sulfated form of dehydroepiandrosterone (DHEA), an androgen precursor produced primarily by the zona reticularis of the adrenal cortex (80--90\%), with a minor contribution from ovarian theca cells in women and testicular Leydig cells in men. DHEA-S is the most abundant circulating steroid hormone in humans. DHEA is synthesized from 17-hydroxypregnenolone via the delta-5 pathway by CYP17A1 (17,20-lyase activity) \cite{auchus2018androgen}. DHEA is then sulfated by DHEA sulfotransferase (SULT2A1) to form DHEA-S, which serves as a storage and transport form \cite{rege2018dhea}. DHEA-S has a much longer plasma half-life (7--10 hours) compared to DHEA (15--30 minutes), exhibits minimal diurnal variation, and is largely protein-bound (albumin). Because of its exclusively adrenal origin, DHEA-S is an excellent marker of adrenal androgen production. Its measurement is clinically useful for: evaluating hyperandrogenism in women (hirsutism, acne, androgenic alopecia, virilization), diagnosing adrenal tumors (adenoma, carcinoma) as opposed to ovarian causes of androgen excess, assessing adrenal reserve in suspected adrenal insufficiency, and monitoring congenital adrenal hyperplasia. DHEA-S declines steadily with age after peaking in the third decade (adrenopause), which is distinct from the gonadal menopause.

\section*{Alternative Names}
\begin{itemize}
  \item DHEA-S
  \item Dehydroepiandrosterone Sulfate
  \item DHEA Sulfate
  \item DHEAS
  \item Serum DHEA-Sulfate
  \item Androstenolone Sulfate
  \item Prasterone Sulfate
\end{itemize}
\section*{Principle}
DHEA-S is measured using competitive immunoassays on automated analyzers. Because DHEA-S circulates at relatively high concentrations (microgram/dL range, approximately 1000 times higher than DHEA), the assay has high analytical precision. An anti-DHEA-S antibody competes with DHEA-S in the patient sample for binding to labeled DHEA-S (chemiluminescent tracer). A displacement agent releases DHEA-S from albumin. The chemiluminescent signal is inversely proportional to DHEA-S concentration. Common platforms: Roche Cobas ECLIA, Abbott Architect CMIA, Siemens Atellica/ADVIA Centaur CLIA, Beckman Coulter DxI. DHEA-S immunoassays have limited cross-reactivity with DHEA (1--5\%), testosterone, androstenedione, and other androgens. The gold standard reference method is liquid chromatography-tandem mass spectrometry (LC-MS/MS), which provides higher specificity, particularly for low DHEA-S levels (as in adrenal insufficiency, children, or elderly), but is not widely available for routine clinical use. DHEA-S measurement does NOT require pre-analytical chilling or special stabilization; it is stable at room temperature for 24 hours and at 2--8 degrees C for 7 days.

\section*{History}
DHEA was first isolated from human urine in 1934 by Adolf Butenandt and colleagues. DHEA-S was identified as the major circulating androgen conjugate in the 1960s by Baulieu and colleagues \cite{baulieu1965adrenalsecreted}. The radioimmunoassay for DHEA-S was developed in the 1970s. The concept of adrenarche (the pubertal increase in adrenal androgen production, particularly DHEA and DHEA-S, between ages 6--8 years, preceding gonadarche) was described by Albright in the 1940s and further characterized by Parker and colleagues. The age-related decline in DHEA-S (adrenopause, starting around age 30) was characterized in large epidemiological studies (the Massachusetts Male Aging Study and others). The use of DHEA-S as a tumor marker for adrenal carcinoma was established in the 1980s--1990s. The role of DHEA-S in PCOS (as a marker of adrenal vs. ovarian hyperandrogenism) was established in the 1990s, with approximately 25--35\% of PCOS women having elevated DHEA-S. DHEA supplementation for "anti-aging" became popular in the late 1990s after the DHEA and Aging book by Baulieu, though randomized trials have not demonstrated consistent benefit and exogenous DHEA (which is converted to DHEA-S) can interfere with test interpretation.

\section*{Why This Test Is Ordered}
\begin{itemize}
  \item Evaluation of hyperandrogenism in women presenting with hirsutism (Ferriman-Gallwey score >8), severe inflammatory acne unresponsive to topical therapy, androgenic alopecia, oligomenorrhea/amenorrhea, or virilization (clitoromegaly, deepening voice, temporal recession, increased muscle mass)
  \item Differentiation of adrenal from ovarian sources of hyperandrogenism --- elevated DHEA-S strongly suggests adrenal origin (adrenal tumor, congenital adrenal hyperplasia); elevated testosterone with normal DHEA-S suggests ovarian origin (PCOS, ovarian hyperthecosis)
  \item Diagnosis and monitoring of adrenal carcinoma --- DHEA-S and DHEA-S are often markedly elevated (>700 mcg/dL, often >1000 mcg/dL) and serve as tumor markers for adrenocortical carcinoma \cite{arlt2011urine}
  \item Screening for congenital adrenal hyperplasia (21-hydroxylase deficiency, 11-beta-hydroxylase deficiency) --- elevated DHEA-S along with 17-hydroxyprogesterone
  \item Investigation of precocious pubarche (early pubic hair development before age 8 in girls and 9 in boys) to distinguish premature adrenarche (mildly elevated DHEA-S for age) from CAH or adrenal tumor
  \item Assessment of adrenal androgen status in suspected adrenal insufficiency --- low DHEA-S (age- and sex-adjusted) supports adrenal androgen deficiency
  \item Monitoring glucocorticoid therapy in congenital adrenal hyperplasia --- DHEA-S suppression indicates adequate adrenal androgen control; target is age-appropriate mid-normal range
  \item Evaluation of suspected exogenous androgen/DHEA abuse (athletic doping) --- elevated DHEA-S with low gonadotropins and low testosterone in males suggests exogenous DHEA ingestion
  \item Assessment of adrenal function in postmenopausal women with symptoms of androgen deficiency (controversial --- insufficient evidence for routine use)
\end{itemize}

\section*{Primary vs Secondary Test}
DHEA-S is an adjunctive (secondary) test in the evaluation of hyperandrogenism. The primary tests for hyperandrogenism are total and free testosterone. DHEA-S is obtained when adrenal source needs to be excluded or confirmed. In PCOS, approximately 25--35\% of patients have elevated DHEA-S, and it is one of the diagnostic criteria in the Rotterdam criteria \cite{teede2018recommendations}. As a tumor marker for adrenal carcinoma, DHEA-S is a primary marker alongside other adrenal steroids.

\section*{How This Test Is Performed}
A healthcare professional draws a blood sample from a vein into a serum separator tube (gold-top) or plain red-top tube. Approximately 3--5 mL of blood is collected. The specimen does not require special handling (no need for chilling or freezing). DHEA-S is stable at room temperature for 24 hours, at 2--8 degrees C for 7 days, and long-term at -20 degrees C \cite{rifai2018tietz}. The draw can be performed at any time of day because DHEA-S does not exhibit significant circadian variation (unlike cortisol and ACTH). However, for consistency in serial monitoring, drawing at approximately the same time is recommended. In menstruating women, DHEA-S does not significantly vary across the menstrual cycle, so timing within the cycle is not critical.

\section*{What Preparation Is Needed}
No fasting is required. DHEA-S does not require special dietary preparation. Patients taking DHEA as an over-the-counter supplement (often sold for "anti-aging," athletic performance, or well-being) must discontinue the supplement for at least 3--5 days before testing because exogenous DHEA is rapidly converted to DHEA-S and will produce falsely elevated results. Oral contraceptives and estrogen therapy decrease DHEA-S by 20--50\% (estrogens increase SHBG, and decrease adrenal androgen production, though the exact mechanism is multifactorial). Glucocorticoids (prednisone >5 mg/day, dexamethasone) suppress ACTH and thereby decrease adrenal androgen production including DHEA-S. High-dose biotin (>5 mg/day) should be discontinued 48 hours before testing because of interference with streptavidin-biotin immunoassay platforms. Insulin, hyperinsulinemia, and obesity are associated with lower DHEA-S levels for unclear reasons.

\section*{Contraindications and Precautions}
\begin{itemize}
  \item The major pre-analytical consideration is the marked age and sex dependence of DHEA-S levels, which MUST be interpreted using age- and sex-specific reference ranges. DHEA-S peaks at ages 20--30 (males 280--640 mcg/dL; females 130--400 mcg/dL), then declines progressively by approximately 2--3\% per year (adrenopause). By age 70--80, DHEA-S levels are approximately 10--20\% of peak values. Without age-adjusted reference ranges, DHEA-S values cannot be correctly interpreted. DHEA-S is falsely decreased by: oral contraceptives and estrogen therapy (20--50\% decrease), glucocorticoid therapy (ACTH suppression), adrenal insufficiency (primary or secondary), critical illness, and late pregnancy (DHEA-S used by the fetal-placental unit for estrogen synthesis; maternal DHEA-S drops by 50\% by third trimester). DHEA-S is elevated in: congenital adrenal hyperplasia, adrenal adenoma/carcinoma, premature adrenarche, polycystic ovary syndrome (PCOS, 25--35\% of patients), insulin resistance paradoxically (though obesity per se is associated with lower DHEA-S), and exogenous DHEA supplementation. Heterophile antibodies and HAMA can interfere with immunoassays. Cross-reactivity with other sulfated steroids (androstenedione sulfate, testosterone sulfate) can produce minor interference but is not clinically significant at the high concentrations measured.
\end{itemize}
\section*{Risks}
Minimal risk. Slight pain, bruising, or hematoma at venipuncture site. Rarely: vasovagal syncope, infection, or phlebitis. No radiation exposure.

\section*{What Happens After the Test}
Results are available within 1--4 hours on automated immunoassay platforms. Interpretation depends on clinical context and age/sex: DHEA-S elevated >95th percentile for age and sex requires further evaluation for adrenal source. In women with hyperandrogenism, an elevated DHEA-S >400 mcg/dL should prompt adrenal imaging (CT or MRI) to evaluate for adrenal adenoma or carcinoma, and measurement of 17-hydroxyprogesterone to exclude late-onset (non-classical) congenital adrenal hyperplasia. Adrenal carcinoma should be excluded when DHEA-S exceeds 700--1000 mcg/dL, especially with rapid-onset virilization. Moderately elevated DHEA-S (200--400 mcg/dL) in hyperandrogenic women is most commonly due to PCOS with an adrenal component. In patients with suspected adrenal insufficiency, a low DHEA-S for age and sex supports the diagnosis, though cortisol and ACTH are the diagnostic tests. Patients with confirmed congenital adrenal hyperplasia and elevated DHEA-S require optimization of glucocorticoid therapy.

\section*{Understanding Results}

\subsection*{Below Normal}
\begin{itemize}
  \item \textbf{Adrenal insufficiency (primary or secondary):} DHEA-S is typically low in both forms. In primary adrenal insufficiency (Addison disease), destruction of the zona reticularis causes loss of adrenal androgen production, and DHEA-S is often undetectable or very low. In women, DHEA replacement (25--50 mg/day) may improve quality of life, libido, and mood (per Endocrine Society guidelines), though evidence is mixed \cite{wierman2014androgen}.
  \item \textbf{Glucocorticoid therapy:} Exogenous glucocorticoids suppress ACTH and thereby suppress adrenal androgen production, including DHEA-S. This occurs with prednisone equivalent >5 mg/day for >3 weeks.
  \item \textbf{Aging (adrenopause):} Physiologic decline in DHEA-S production from the adrenal zona reticularis that begins at approximately age 30 and declines by 2--3\% per year. By ages 70--80, DHEA-S levels are 10--20\% of peak young adult values. This is a normal aging process, not a disease.
  \item \textbf{Hypopituitarism and Sheehan syndrome:} Loss of ACTH secretion leads to reduced adrenal androgen production and low DHEA-S.
  \item \textbf{Chronic illness:} Severe chronic disease of any type (cancer, heart failure, renal failure, liver disease, HIV/AIDS) can suppress DHEA-S production through unclear mechanisms, possibly involving cytokines and decreased ACTH drive.
  \item \textbf{Oral contraceptive use and estrogen therapy:} Decrease DHEA-S by 20--50\%.
  \item \textbf{Anorexia nervosa and severe malnutrition:} Profoundly low DHEA-S with preservation or mild decrease in cortisol.
  \item \textbf{Late pregnancy (third trimester):} Maternal DHEA-S declines by 50\% because the fetal-placental unit utilizes DHEA-S as a substrate for estrogen synthesis (placental sulfatase cleaves DHEA-S to DHEA, which is aromatized to estrogens).
\end{itemize}

\subsection*{Above Normal}
\begin{itemize}
  \item \textbf{Adrenocortical carcinoma:} Markedly elevated DHEA-S (>700 mcg/dL, often >1000 mcg/dL) is characteristic. DHEA-S and other androgen precursors are produced by the tumor (which is often hormonally active). Rapid-onset hirsutism, virilization, and cushingoid features may occur simultaneously. Tumor is typically >6 cm, heterogeneous, and irregular on CT. DHEA-S serves as a tumor marker for monitoring after resection.
  \item \textbf{Adrenal adenoma:} Mild-to-moderate elevation of DHEA-S, typically <700 mcg/dL. Usually non-functional (incidentaloma) or producing mild androgen excess. Adrenal CT shows homogeneous, well-circumscribed mass <4 cm with lipid-rich attenuation (<10 HU).
  \item \textbf{Congenital adrenal hyperplasia (21-hydroxylase deficiency):} DHEA-S is markedly elevated in the classical (severe) form, along with elevated 17-hydroxyprogesterone and androstenedione. In non-classical (late-onset) CAH, DHEA-S may be mildly elevated or normal. 17-hydroxyprogesterone after ACTH stimulation is diagnostic (>1000--1500 ng/dL).
  \item \textbf{Polycystic ovary syndrome (PCOS):} Approximately 25--35\% of PCOS women have elevated DHEA-S (adrenal component of PCOS hyperandrogenism). DHEA-S typically 200--400 mcg/dL, rarely >400 mcg/dL. Testosterone is usually elevated as well. DHEA-S elevation identifies the adrenal component of the syndrome.
  \item \textbf{Premature adrenarche:} Early pubic hair development (before age 8 in girls, 9 in boys) with mildly elevated DHEA-S for age but appropriate for Tanner stage. This is a benign, non-progressive condition in most cases but requires exclusion of CAH and adrenal tumor.
  \item \textbf{Exogenous DHEA supplementation:} DHEA is available over-the-counter as a dietary supplement (prasterone). Ingestion results in elevated DHEA-S, which may confound clinical testing. Patients should discontinue supplementation 3--5 days before testing.
  \item \textbf{Adrenal androgen-secreting tumors other than carcinoma:} Rare virilizing adenomas can produce primarily DHEA-S and DHEA with or without cortisol excess. Suspect when DHEA-S >400 mcg/dL without other explanation.
  \item \textbf{Cushing disease:} Moderate DHEA-S elevation can occur due to ACTH-driven adrenal androgen production alongside hypercortisolism, particularly in ectopic ACTH syndrome.
\end{itemize}

\section*{Normal Results}
\begin{itemize}
  \item \textbf{DHEA-S, Males 18--29 years:} 280--640 mcg/dL (7.6--17.4 micromol/L; Roche ECLIA)
  \item \textbf{DHEA-S, Females 18--29 years:} 130--400 mcg/dL (3.5--10.9 micromol/L)
  \item \textbf{DHEA-S, Males 30--39 years:} 150--450 mcg/dL
  \item \textbf{DHEA-S, Females 30--39 years:} 80--330 mcg/dL
  \item \textbf{DHEA-S, Males 40--49 years:} 120--380 mcg/dL
  \item \textbf{DHEA-S, Females 40--49 years:} 60--280 mcg/dL
  \item \textbf{DHEA-S, Males 50--59 years:} 80--300 mcg/dL
  \item \textbf{DHEA-S, Females 50--59 years:} 40--220 mcg/dL
  \item \textbf{DHEA-S, Males 60--69 years:} 50--220 mcg/dL
  \item \textbf{DHEA-S, Females 60--69 years:} 25--150 mcg/dL
  \item \textbf{DHEA-S, Males >70 years:} 25--150 mcg/dL
  \item \textbf{DHEA-S, Females >70 years:} 15--100 mcg/dL
  \item \textbf{DHEA-S, Children <10 years:} <100 mcg/dL (pre-pubertal); increases during adrenarche (age 6--8) and peaks at age 20--30
  \item \textbf{DHEA-S, Newborns (first week):} High neonatal DHEA-S from fetal zone of adrenal cortex (levels 100--500 mcg/dL) which declines by 1--2 months
\end{itemize}

\section*{Follow-up Tests}
\begin{itemize}
  \item \textbf{Total and free testosterone:} Primary tests for hyperandrogenism. Elevated free testosterone is the hallmark of hyperandrogenism. Co-measurement with DHEA-S determines adrenal vs. ovarian source.
  \item \textbf{17-hydroxyprogesterone (17-OHP):} Screening test for 21-hydroxylase deficiency (congenital adrenal hyperplasia) in patients with elevated DHEA-S. Basal 17-OHP >200 ng/dL (follicular phase) merits ACTH stimulation testing. 17-OHP >1000--1500 ng/dL after ACTH is diagnostic for non-classical CAH.
  \item \textbf{Androstenedione:} Additional adrenal androgen marker that is often elevated in CAH and adrenal tumors. Useful in PCOS evaluation as a secondary marker.
  \item \textbf{ACTH (plasma):} Distinguishes ACTH-dependent (pituitary, ectopic) from ACTH-independent (adrenal tumor) causes of elevated DHEA-S.
  \item \textbf{ACTH stimulation test (250 mcg cosyntropin):} Evaluates adrenal androgen reserve. 17-OHP >1500 ng/dL at 60 minutes is diagnostic for 21-hydroxylase deficiency. Used for non-classical CAH screening.
  \item \textbf{Adrenal CT (non-contrast with washout):} Imaging for suspected adrenal adenoma or carcinoma when DHEA-S is significantly elevated or virilization is present.
  \item \textbf{Pelvic ultrasound:} Investigates ovarian source of hyperandrogenism (polycystic ovarian morphology --- >20 follicles per ovary or ovarian volume >10 mL). Used in PCOS diagnosis per Rotterdam criteria.
  \item \textbf{Dexamethasone suppression test:} Suppression of DHEA-S by >50\% after 2--4 days of dexamethasone suggests ACTH-dependent androgen excess (functional adrenal hyperandrogenism); lack of suppression suggests autonomous adrenal tumor.
  \item \textbf{LH, FSH, and SHBG:} Part of the full PCOS evaluation. LH:FSH ratio >2--3 is commonly seen in PCOS. SHBG is typically low in PCOS due to hyperinsulinemia, contributing to elevated free testosterone.
\end{itemize}

\section*{References}
\bibliographystyle{unsrtnat}
\bibliography{references}

\clearpage
\section*{Regulatory Status and Authorization Date}
This chapter describes a test category, not a specific commercial product. FDA, CDSCO, EU IVDR, and MHRA approval or registration dates therefore cannot be assigned without the assay manufacturer, product name, instrument, and intended use. For laboratory-developed tests, an individual agency approval date may not exist. See the Preface for the interpretation of regulatory status.

\clearpage
